\documentclass[11pt,a4paper]{article}

\usepackage[utf8]{inputenc}
\usepackage[T1]{fontenc}
\usepackage[italian]{babel}
\usepackage{amsmath}
\usepackage{amssymb}
\usepackage[margin=2.2cm]{geometry}
\usepackage{xcolor}
\usepackage{listings}
\usepackage{enumitem}
\usepackage{booktabs}
\usepackage{longtable}
\usepackage{array}
\usepackage{titlesec}
\usepackage[hidelinks]{hyperref}

\definecolor{codebg}{rgb}{0.96,0.96,0.96}
\definecolor{kw}{rgb}{0.0,0.2,0.6}
\definecolor{cmt}{rgb}{0.3,0.5,0.3}

\lstset{
  basicstyle=\ttfamily\scriptsize,
  backgroundcolor=\color{codebg},
  keywordstyle=\color{kw}\bfseries,
  commentstyle=\color{cmt}\itshape,
  frame=single,
  framesep=4pt,
  breaklines=true,
  showstringspaces=false,
  numbers=left,
  numberstyle=\tiny\color{gray},
  language=C,
  tabsize=4
}

\newcommand{\code}[1]{\texttt{#1}}
\newcommand{\func}[1]{\textbf{\texttt{#1}}}


\titleformat{\section}{\large\bfseries\color{kw}}{\thesection}{0.6em}{}
\titleformat{\subsection}{\normalsize\bfseries}{\thesubsection}{0.6em}{}
\titleformat{\subsubsection}{\normalsize\bfseries\itshape}{\thesubsubsection}{0.6em}{}

\setlist{itemsep=1pt,topsep=2pt}

\title{\textbf{PandOSsh --- Fase 1}\\[2pt]
\large Le strutture dati di base: liste, PCB e ASL\\
\normalsize Guida dettagliata riga per riga, funzione per funzione}
\author{Materiale di studio per l'interrogazione}
\date{A.A. 2025/2026}

\begin{document}
\maketitle
\tableofcontents
\bigskip
\hrule
\bigskip

% ============================================================
\section{Introduzione generale alla Fase 1}
% ============================================================

La \textbf{Phase 1} di PandOSsh non è ancora un sistema operativo: non
esistono scheduler, non esistono eccezioni gestite, non gira nulla in
autonomia. È il livello più basso costruito \emph{sopra} il firmware
dell'emulatore $\mu$RISCV (che fa da BIOS), e fornisce esclusivamente
\textbf{strutture dati di servizio} che i livelli superiori (Nucleus in
Phase~2, Support Level in Phase~3) utilizzeranno per implementare i concetti
di processo, sincronizzazione e scheduling.

Il codice della Fase 1 è interamente contenuto in due moduli:

\begin{itemize}
\item \code{phase1/pcb.c} --- gestione dei \textbf{Process Control Block}
  (PCB): allocazione/deallocazione, code di processi, albero dei processi.
\item \code{phase1/asl.c} --- gestione della \textbf{Active Semaphore List}
  (ASL): la struttura che tiene traccia di quali processi sono bloccati su
  quale semaforo.
\end{itemize}

Entrambi i moduli si appoggiano a una libreria di liste concatenate generica
(\code{headers/listx.h}, un sottoinsieme delle liste del kernel Linux) e alle
definizioni di tipo in \code{headers/types.h}.

\paragraph{Principio cardine: niente allocazione dinamica.}
In tutta la Fase 1 (e in tutto il progetto) non compare mai \code{malloc}.
Il numero massimo di processi (\code{MAXPROC}, 20) e di semafori attivi è
fissato a tempo di compilazione: tutte le strutture sono array statici,
gestiti tramite \emph{free list}. Questo è un requisito tipico dei sistemi
operativi ``didattici in stile Pandos'': un kernel deve avere un
comportamento \textbf{deterministico}, senza frammentazione della memoria e
senza il rischio che un'allocazione fallisca in un momento critico.

\paragraph{Perché liste doppiamente concatenate circolari?}
Sia i PCB sia i semafori attivi sono organizzati con la stessa tecnica: una
\code{struct list\_head} embedded nella struttura dati, gestita con le
funzioni generiche di \code{listx.h}. I vantaggi:
\begin{itemize}
\item inserimento e rimozione in $O(1)$, qualunque sia la posizione (testa,
  coda, elemento intermedio);
\item nessun caso speciale per lista vuota o con un solo elemento (la
  sentinella \code{head} è essa stessa un nodo, quindi il codice che
  manipola \code{next}/\code{prev} è identico in ogni caso);
\item una stessa struttura (es. \code{pcb\_t}) può appartenere
  contemporaneamente a più liste diverse (free list, ready queue, coda di un
  semaforo, lista dei figli...) semplicemente avendo più campi
  \code{list\_head}.
\end{itemize}

% ============================================================
\section{\code{headers/types.h} --- i tipi di dato fondamentali}
% ============================================================

Questo header definisce \emph{tutte} le strutture dati condivise da tutto il
progetto (non solo dalla Fase 1). È importante conoscerle a menadito perché
compaiono ovunque nel resto del kernel.

\begin{lstlisting}
typedef signed int cpu_t;
typedef unsigned int memaddr;

/* Page Table Entry descriptor */
typedef struct pteEntry_t
{
    unsigned int pte_entryHI;
    unsigned int pte_entryLO;
} pteEntry_t;

/* Support level context */
typedef struct context_t
{
    unsigned int stackPtr;
    unsigned int status;
    unsigned int pc;
} context_t;

/* Support level descriptor */
typedef struct support_t
{
    int sup_asid;
    state_t sup_exceptState[2];
    context_t sup_exceptContext[2];
    pteEntry_t sup_privatePgTbl[USERPGTBLSIZE];
    unsigned int sup_stackTLB[500];
    unsigned int sup_stackGen[500];
    struct list_head s_list;
} support_t;

/* Page swap pool information structure type */
typedef struct swap_t
{
    int sw_asid;
    int sw_pageNo;
    pteEntry_t *sw_pte;
} swap_t;

/* process table entry type */
typedef struct pcb_t {
    struct list_head p_list;

    struct pcb_t    *p_parent;
    struct list_head p_child;
    struct list_head p_sib;

    state_t p_s;
    cpu_t   p_time;

    int *p_semAdd;
    support_t *p_supportStruct;

    int p_prio;
    int p_pid;
} pcb_t, *pcb_PTR;

/* semaphore descriptor (SEMD) data structure */
typedef struct semd_t {
    int *s_key;
    struct list_head s_procq;
    struct list_head s_link;
} semd_t, *semd_PTR;
\end{lstlisting}

\subsection{\code{cpu\_t} e \code{memaddr}}
Due semplici \code{typedef} per rendere il codice autodocumentante:
\code{cpu\_t} (int con segno) misura quantità di tempo CPU (letto dal TOD,
il Time-Of-Day clock), \code{memaddr} (unsigned int) rappresenta un
indirizzo di memoria a 32 bit.

\subsection{\code{pteEntry\_t} --- entry di Page Table}
Rappresenta \emph{una} riga della page table usata dalla memoria virtuale
(Fase 3): \code{pte\_entryHI} contiene VPN (Virtual Page Number) e ASID,
\code{pte\_entryLO} contiene PFN (Physical Frame Number) più i bit di
controllo (\code{V} = valida, \code{D} = dirty/scrivibile). Non è usata dalla
Fase 1, ma è dichiarata qui perché fa parte del ``vocabolario'' di tipi di
tutto il progetto.

\subsection{\code{context\_t} --- contesto per il pass-up}
Tripletta \emph{(stack pointer, status, program counter)} che descrive
\emph{dove} il Nucleus deve saltare quando passa un'eccezione a un livello
superiore (Support Level). Viene caricata dall'istruzione \code{LDCXT} in
\code{passUpOrDie} (Fase 2).

\subsection{\code{support\_t} --- la ``Support Structure''}
È la struttura che collega Nucleus e Support Level (dettagliata in Fase 3):
contiene l'ASID del processo, due slot di stato salvato
(\code{sup\_exceptState[2]}, uno per i page fault e uno per le altre
eccezioni), due context di ripartenza (\code{sup\_exceptContext[2]}), la
Page Table privata del processo utente (32 entry) e due stack dedicati agli
handler (\code{sup\_stackTLB}, \code{sup\_stackGen}, 500 word ciascuno). Ogni
\code{pcb\_t} ha un puntatore opzionale a una \code{support\_t}
(\code{p\_supportStruct}): se è \code{NULL} il processo è un processo
``puro kernel'' (Fase 2), altrimenti è una U-proc (Fase 3).

\subsection{\code{swap\_t} --- descrittore di frame nello Swap Pool}
Usato solo in Fase 3: per ogni frame fisico dello Swap Pool ricorda quale
ASID e quale pagina logica vi risiedono, e un puntatore alla PTE
corrispondente (per poterla invalidare rapidamente in caso di sfratto).

\subsection{\code{pcb\_t} --- il Process Control Block}
\label{sub:pcb-t-fase1}
Il cuore della Fase 1. Campo per campo:

\begin{center}
\begin{tabular}{@{}lp{10.3cm}@{}}
\toprule
\textbf{Campo} & \textbf{Significato} \\
\midrule
\code{p\_list} & nodo di lista usato per inserire il PCB in \emph{una} coda alla
  volta: la free list dei PCB liberi, la ready queue, oppure la coda dei
  processi bloccati su un semaforo (\code{s\_procq} in \code{semd\_t}). Un
  PCB non può stare in due di queste code contemporaneamente, quindi
  riutilizzare lo stesso campo per tutte e tre è sicuro e risparmia memoria. \\
\code{p\_parent} & puntatore al PCB del processo padre (\code{NULL} per il
  processo radice). \\
\code{p\_child} & testa della lista dei figli di \emph{questo} processo. \\
\code{p\_sib} & nodo con cui questo PCB si aggancia alla lista \code{p\_child}
  del proprio padre (lista dei ``fratelli''). \\
\code{p\_s} & \code{state\_t}: l'intero stato del processore (registri,
  PC, status, ecc.) salvato/ripristinato a ogni context switch. \\
\code{p\_time} & tempo di CPU accumulato dal processo (usato dalla syscall
  \code{GETTIME}). \\
\code{p\_semAdd} & puntatore intero: se non \code{NULL}, indirizzo del
  semaforo su cui il processo è bloccato (altrimenti il processo è pronto o
  in esecuzione). \\
\code{p\_supportStruct} & puntatore alla \code{support\_t} del processo, o
  \code{NULL} se non ne ha una. \\
\code{p\_prio} & priorità (\code{PROCESS\_PRIO\_LOW}=0 o
  \code{PROCESS\_PRIO\_HIGH}=1) usata dallo scheduler. \\
\code{p\_pid} & identificatore del processo, assegnato in ordine crescente
  da \code{allocPcb}. \\
\bottomrule
\end{tabular}
\end{center}

\textbf{Da notare}: \code{p\_list} è usato per \emph{tre} scopi diversi in
tempi diversi (free list / ready queue / coda semaforo) ma mai
simultaneamente: è un classico esempio di riuso di un campo di collegamento
generico offerto dalla libreria di liste intrusive.

\subsection{\code{semd\_t} --- descrittore di semaforo attivo}
\begin{itemize}
\item \code{s\_key}: puntatore che identifica \emph{univocamente} il
  semaforo (in pratica l'indirizzo della variabile intera che rappresenta il
  semaforo, es. un elemento dell'array \code{devSems}).
\item \code{s\_procq}: testa della coda dei PCB bloccati su questo semaforo
  (i nodi agganciati sono i campi \code{p\_list} dei PCB).
\item \code{s\_link}: nodo con cui il descrittore si aggancia alla ASL
  (\code{semd\_h}) oppure alla free list dei descrittori
  (\code{semdFree\_h}).
\end{itemize}

% ============================================================
\section{Tipi forniti dalla SDK \texorpdfstring{$\mu$}{u}RISCV (\code{uriscv/types.h})}
% ============================================================

\code{headers/types.h} inizia con \code{\#include <uriscv/types.h>}: un
header di sistema (in \code{/usr/include/uriscv/}, \textbf{non} scritto dal
progetto) che definisce i tipi legati direttamente all'hardware emulato.
Questi tipi compaiono in \emph{ogni} listato dei tre documenti (\code{p\_s},
\code{savedState}, \code{exState}\ldots) ma non sono mai dichiarati nei
sorgenti del progetto: vale la pena vederli una volta per tutte qui, perché
altrimenti campi come \code{->reg\_a0} o \code{->cause} restano ``magici''.

\subsection{\code{state\_t} --- lo stato del processore}
\begin{lstlisting}
#define STATE_GPR_LEN 32

typedef struct state {
    unsigned int entry_hi;
    unsigned int cause;
    unsigned int status;
    unsigned int pc_epc;
    unsigned int mie;
    unsigned int gpr[STATE_GPR_LEN];
} state_t;

/* alias per i 32 registri general-purpose (convenzione RISC-V) */
#define reg_sp gpr[2]
#define reg_a0 gpr[24]
#define reg_a1 gpr[25]
#define reg_a2 gpr[26]
#define reg_a3 gpr[27]
/* ... reg_ra, reg_gp, reg_t0..t6, reg_s0..s11, reg_a4..a7 ... */
\end{lstlisting}
\`E la struttura \textbf{più importante di tutto il progetto}: rappresenta
l'intero stato del processore in un dato istante, ed è ciò che viene
salvato/ripristinato a ogni context switch, eccezione o interrupt.
\begin{itemize}
\item \code{entry\_hi}: registro TLB con VPN (indirizzo virtuale) e ASID
  --- rilevante nei page fault (Fase 3).
\item \code{cause}: il registro che spiega \emph{perché} è avvenuta
  l'eccezione (bit 31 = interrupt; altrimenti \code{ExcCode} nei bit
  bassi). \`E il primo campo letto da \code{exceptionHandler}.
\item \code{status}: i bit di modalità e abilitazione interrupt
  (\code{MPP} = kernel/user mode, \code{MIE}/\code{MPIE} = interrupt
  correnti/precedenti abilitati).
\item \code{pc\_epc}: il program counter, cioè l'indirizzo dell'istruzione
  a cui riprendere l'esecuzione.
\item \code{mie}: maschera di quali sorgenti di interrupt sono abilitate.
\item \code{gpr[32]}: i 32 registri general-purpose della RISC-V
  (\code{x0}--\code{x31}). Il progetto non li indicizza mai numericamente,
  ma tramite le macro \code{reg\_sp}, \code{reg\_a0}--\code{reg\_a3}
  (usate per leggere il numero e gli argomenti di una \code{SYSCALL},
  esattamente come \code{a0}--\code{a3} nella convenzione di chiamata
  RISC-V), ecc.
\end{itemize}
Nel PCB, il campo \code{p\_s} (Sezione~\ref{sub:pcb-t-fase1}) è di tipo \code{state\_t}: è
\emph{la} copia persistente dello stato di un processo mentre non sta
girando. Quando un'eccezione avviene, il \textbf{firmware} (non il
Nucleus) salva automaticamente lo stato corrente in un'area fissa di
memoria (\code{BIOSDATAPAGE}, si veda il documento di Fase 2): è da lì che
\code{exceptionHandler} legge \code{savedState}.

\subsection{\code{passupvector\_t} --- il Pass Up Vector}
\begin{lstlisting}
typedef struct passupvector {
    unsigned int tlb_refill_handler;
    unsigned int tlb_refill_stackPtr;
    unsigned int exception_handler;
    unsigned int exception_stackPtr;
} passupvector_t;
\end{lstlisting}
Quattro interi che il Nucleus scrive \textbf{una sola volta}, in
\code{initial.c::main} (si veda il documento di Fase 2), a un indirizzo
fisso (\code{PASSUPVECTOR}, \code{0x0FFFF900}) letto dal firmware: dicono
all'hardware dove saltare (\code{\_handler}) e quale stack usare
(\code{\_stackPtr}) rispettivamente per un evento di TLB-Refill e per
qualunque altra eccezione. \`E il meccanismo con cui il Nucleus si
``registra'' presso il firmware come gestore di tutte le eccezioni.

\subsection{\code{dtpreg\_t} e \code{termreg\_t} --- registri dei device}
\begin{lstlisting}
/* Device "semplici": disk, flash, printer */
typedef struct dtpreg {
    unsigned int status;
    unsigned int command;
    unsigned int data0;
    unsigned int data1;
} dtpreg_t;

/* Terminali: due sotto-canali indipendenti (ricezione e trasmissione) */
typedef struct termreg {
    unsigned int recv_status;
    unsigned int recv_command;
    unsigned int transm_status;
    unsigned int transm_command;
} termreg_t;
\end{lstlisting}
Due modi di interpretare il blocco di registri hardware di un device
(l'indirizzo di partenza si calcola con la macro \code{DEV\_REG\_ADDR(line,
dev)}, anch'essa fornita dalla SDK in \code{uriscv/arch.h}). I device
``semplici'' (disk, flash, printer) hanno un'unica coppia
status/comando (più due registri dati, \code{data0}/\code{data1}, usati
per il DMA: si veda \code{flashOperation} in Fase 3, che scrive l'indirizzo
del frame in \code{data0}). I terminali, avendo due canali fisicamente
indipendenti, raddoppiano la coppia: \code{recv\_*} per la ricezione,
\code{transm\_*} per la trasmissione --- esattamente la distinzione già
introdotta a proposito di \code{TERM\_TX\_SEM}/\code{TERM\_RX\_SEM}
(Fase 2) e usata da \code{writeTerminal}/\code{readTerminal} (Fase 3).
Si noti che \code{interrupts.c} (Fase 2) non usa questi due \code{typedef}
per nome, ma ottiene lo stesso effetto indicizzando manualmente un array
di \code{unsigned int} agli stessi offset (si vedano le macro
\code{DEV\_STATUS}/\code{DEV\_COMMAND}/\code{TERM\_RECV\_STATUS}\ldots nel
documento di Fase 2): risultato identico, stile diverso.

% ============================================================
\section{\code{headers/listx.h} --- liste concatenate generiche}
% ============================================================

Questo file è un sottoinsieme (semplificato) delle liste doppiamente
concatenate circolari del kernel Linux (\code{include/linux/list.h}). È la
base tecnica su cui si costruiscono \emph{tutte} le code del progetto: ready
queue, code dei semafori, lista dei figli, free list dei PCB e dei
\code{semd\_t}.

\subsection{L'idea: liste ``intrusive''}
Invece di una lista generica che contiene puntatori a dati arbitrari
(come una \code{List<T>} ad alto livello), qui la lista è \emph{intrusiva}:
un campo \code{struct list\_head} viene inserito \emph{dentro} la struttura
dati che si vuole collegare (es. \code{p\_list} dentro \code{pcb\_t}). La
lista collega direttamente questi campi tra loro; per risalire alla struttura
dati che li contiene si usa la macro \func{container\_of}.

\begin{lstlisting}
struct list_head {
    struct list_head *next, *prev;
};
\end{lstlisting}
Una lista è quindi semplicemente un nodo con due puntatori. Una lista
\emph{vuota} è un nodo il cui \code{next} e \code{prev} puntano \emph{a se
stesso}: questo è il trucco che elimina tutti i casi speciali.

\subsection{\func{offsetof} e \func{container\_of}}
\begin{lstlisting}
#define offsetof(TYPE, MEMBER) ((size_tt)(&((TYPE *)0)->MEMBER))

#define container_of(ptr, type, member)                                    \
    ({                                                                     \
        const typeof(((type *)0)->member) *__mptr = (ptr);                 \
        (type *)((char *)__mptr - offsetof(type, member));                 \
    })
\end{lstlisting}
\begin{itemize}
\item \func{offsetof(TYPE, MEMBER)} calcola, senza mai istanziare
  davvero una \code{TYPE}, la distanza in byte dall'inizio della struttura al
  campo \code{MEMBER}: finge che una \code{TYPE} sia allocata all'indirizzo
  \code{0} e legge l'indirizzo (che numericamente coincide con l'offset) del
  campo richiesto.
\item \func{container\_of(ptr, type, member)} fa l'operazione inversa
  di ``prendere l'indirizzo di un campo e tornare al contenitore'': dato un
  puntatore \code{ptr} a un campo \code{member} di tipo \code{type}, sottrae
  l'offset del campo dall'indirizzo per ottenere il puntatore alla struttura
  intera. È così che, ad esempio, da un \code{struct list\_head *} estratto
  dalla ready queue si risale al \code{pcb\_t*} che lo contiene:
  \code{container\_of(pos, pcb\_t, p\_list)}.
\end{itemize}

\subsection{Inizializzazione}
\begin{lstlisting}
#define LIST_HEAD_INIT(name) { &(name), &(name) }
#define LIST_HEAD(name) struct list_head name = LIST_HEAD_INIT(name)

static inline void INIT_LIST_HEAD(struct list_head *list) {
    list->next = list;
    list->prev = list;
}
\end{lstlisting}
\code{LIST\_HEAD\_INIT}/\code{LIST\_HEAD} sono macro usate per dichiarare e
inizializzare, in un colpo solo, una variabile lista vuota a livello
statico. \func{INIT\_LIST\_HEAD} è la funzione (inline) equivalente da
chiamare a runtime su una lista già dichiarata: imposta entrambi i puntatori
sulla lista stessa, rendendola una lista vuota valida.

\subsection{Inserimento}
\begin{lstlisting}
static inline void __list_add(struct list_head *new, struct list_head *prev, struct list_head *next) {
    next->prev = new;
    new->next  = next;
    new->prev  = prev;
    prev->next = new;
}

static inline void list_add(struct list_head *new, struct list_head *head) {
    __list_add(new, head, head->next);
}

static inline void list_add_tail(struct list_head *new, struct list_head *head) {
    __list_add(new, head->prev, head);
}
\end{lstlisting}
\func{\_\_list\_add} è la primitiva di base: inserisce \code{new} \emph{tra}
due nodi già collegati (\code{prev} e \code{next}), aggiornando i 4
puntatori coinvolti. È $O(1)$ e funziona identicamente sia che la lista sia
vuota, con un elemento o con molti.
\begin{itemize}
\item \func{list\_add(new, head)}: inserisce \code{new} \emph{subito dopo}
  \code{head}, cioè in \textbf{testa} alla lista (\code{prev=head},
  \code{next=head->next}).
\item \func{list\_add\_tail(new, head)}: inserisce \code{new} \emph{subito
  prima} di \code{head}, cioè in \textbf{coda} (\code{prev=head->prev},
  \code{next=head}). È la funzione più usata nel progetto (free list, code
  FIFO, inserimento dei figli).
\end{itemize}

\subsection{Rimozione}
\begin{lstlisting}
static inline void __list_del(struct list_head *prev, struct list_head *next) {
    next->prev = prev;
    prev->next = next;
}

static inline void list_del(struct list_head *entry) {
    __list_del(entry->prev, entry->next);
    entry->next = entry;
    entry->prev = entry;
}
\end{lstlisting}
\func{\_\_list\_del} ricollega direttamente \code{prev} e \code{next}
``scavalcando'' il nodo che sta in mezzo. \func{list\_del(entry)} la usa per
staccare \code{entry} dalla lista in cui si trova, e poi
\textbf{reimposta \code{entry} come lista vuota} (punta a se stesso): questo
evita puntatori pendenti (dangling) se qualcuno riusa quel nodo per
errore prima di reinserirlo altrove.

\subsection{Test e navigazione}
\begin{lstlisting}
static inline int list_is_last(const struct list_head *list, const struct list_head *head) {
    return list->next == head;
}

static inline int list_empty(const struct list_head *head) {
    return head->next == head;
}

static inline struct list_head *list_next(const struct list_head *current) {
    if (list_empty(current)) return NULL;
    else return current->next;
}

static inline struct list_head *list_prev(const struct list_head *current) {
    if (list_empty(current)) return NULL;
    else return current->prev;
}
\end{lstlisting}
\func{list\_empty} sfrutta esattamente l'invariante di inizializzazione: una
lista vuota ha \code{head->next == head}. \func{list\_is\_last} controlla se
\code{list} è l'ultimo nodo prima della sentinella \code{head}.
\func{list\_next}/\func{list\_prev} sono getter di comodo (usati poco nel
progetto, che preferisce le macro di iterazione sottostanti).

\subsection{Macro di iterazione}
\begin{lstlisting}
#define list_for_each(pos, head) \
    for (pos = (head)->next; pos != (head); pos = pos->next)

#define list_for_each_prev(pos, head) \
    for (pos = (head)->prev; pos != (head); pos = pos->prev)

#define list_for_each_entry(pos, head, member)                              \
    for (pos = container_of((head)->next, typeof(*pos), member);            \
         &pos->member != (head);                                            \
         pos = container_of(pos->member.next, typeof(*pos), member))

#define list_for_each_entry_reverse(pos, head, member)                      \
    for (pos = container_of((head)->prev, typeof(*pos), member);            \
         &pos->member != (head);                                            \
         pos = container_of(pos->member.prev, typeof(*pos), member))
\end{lstlisting}
\begin{itemize}
\item \func{list\_for\_each(pos, head)}: itera sui \emph{nodi
  list\_head} grezzi, partendo dal primo elemento reale
  (\code{head->next}) e fermandosi quando si torna alla sentinella. Usata
  in \code{asl.c} e \code{pcb.c} quando serve confrontare \code{pos} con
  l'indirizzo di un campo specifico (es. \code{outProcQ}, \code{outBlocked}).
\item \func{list\_for\_each\_entry(pos, head, member)}: variante di
  comodo che restituisce direttamente il puntatore alla \emph{struttura
  contenitore} (tramite \code{container\_of} a ogni iterazione), evitando di
  dover fare la conversione manualmente nel corpo del ciclo. È quella usata
  in \code{asl.c} per scorrere \code{semd\_h} (\code{semd\_t *s}).
\item Le due varianti \code{\_reverse} percorrono la lista all'indietro
  partendo da \code{head->prev}; non sono usate nel codice del progetto ma
  sono disponibili nella libreria.
\end{itemize}

% ============================================================
\section{\code{phase1/pcb.c} --- Process Control Block}
% ============================================================

\subsection{Il file header \code{phase1/headers/pcb.h}}
\begin{lstlisting}
// SPDX-FileCopyrightText: 2022 Luca Bassi, Gaia Clerici, Mirco Dondi, Fabio Gaiba
//
// SPDX-License-Identifier: GPL-3.0-or-later

#ifndef PCB_H_INCLUDED
#define PCB_H_INCLUDED

#include "../../headers/listx.h"
#include "../../headers/types.h"

void initPcbs();
void freePcb(pcb_t* p);
pcb_t* allocPcb();

void mkEmptyProcQ(struct list_head* head);
int emptyProcQ(struct list_head* head);
void insertProcQ(struct list_head* head, pcb_t* p);
pcb_t* headProcQ(struct list_head* head);
pcb_t* removeProcQ(struct list_head* head);
pcb_t* outProcQ(struct list_head* head, pcb_t* p);

int emptyChild(pcb_t* p);
void insertChild(pcb_t* prnt, pcb_t* p);
pcb_t* removeChild(pcb_t* p);
pcb_t* outChild(pcb_t* p);

#endif
\end{lstlisting}
\code{pcb.h} è l'\textbf{interfaccia pubblica} del modulo: dichiara i
prototipi delle 13 funzioni analizzate più sotto (nel file originale ogni
prototipo è preceduto da un commento su una riga che ne riassume lo scopo
--- lo stesso testo usato come base per le spiegazioni di questa sezione),
include \code{listx.h} (per \code{struct list\_head}) e \code{types.h}
(per \code{pcb\_t}), e non definisce \emph{nulla}: nessuna variabile,
nessun corpo di funzione. \`E questa separazione header/implementazione
che permette a \code{phase2} e \code{phase3} di usare i PCB
(\code{\#include "../phase1/headers/pcb.h"}) senza conoscere né poter
accedere a \code{pcbFree\_h}, \code{pcbFree\_table} o \code{next\_pid}
(dichiarati \code{static} in \code{pcb.c}, quindi invisibili fuori da quel
file): l'unico modo per manipolare i PCB dall'esterno è tramite queste 13
funzioni.

\`E interessante notare che la testata \textbf{SPDX} attribuisce questo
file (e il gemello \code{asl.h}) a un altro gruppo di autori (2022): è
l'\textbf{interfaccia data dal corso} (il ``contratto'' che l'implementazione
deve rispettare, ereditato da un precedente progetto PandOS didattico
open-source), mentre \code{pcb.c} e \code{asl.c} --- le implementazioni
vere e proprie analizzate in questo documento --- sono il lavoro originale
di questo team. Distinguere le due cose è utile se l'esaminatore chiede
``questo file l'avete scritto voi?''.

\subsection{Variabili statiche del modulo}
\begin{lstlisting}
static struct list_head pcbFree_h;
static pcb_t pcbFree_table[MAXPROC];
static int next_pid = 1;
\end{lstlisting}
\begin{itemize}
\item \code{pcbFree\_table[MAXPROC]}: l'unico array di PCB dell'intero
  sistema (20 elementi, \code{MAXPROC=20} in \code{const.h}). \emph{Ogni}
  processo che esisterà mai durante l'esecuzione userà uno di questi 20
  slot: non c'è altro modo per ottenere un PCB.
\item \code{pcbFree\_h}: testa della free list che collega, tramite il
  campo \code{p\_list}, tutti gli slot di \code{pcbFree\_table} attualmente
  \emph{non} in uso.
\item \code{next\_pid}: contatore globale monotono crescente, usato per
  assegnare un PID univoco a ogni PCB allocato. Parte da 1 (0 è riservato
  per rappresentare ``nessun PID''/errore in alcune chiamate, es.
  \code{GETPROCESSID}). \textbf{Nota:} i PID non vengono mai riusati anche se
  un processo termina, quindi crescono indefinitamente entro i limiti di un
  \code{int} --- comportamento accettabile perché il numero di processi
  creabili in una sessione è comunque limitato dalla dimensione della RAM
  simulata e dalla durata delle demo.
\end{itemize}

\subsection{\func{initPcbs} --- inizializzazione}
\begin{lstlisting}
void initPcbs() {
    int i;
    INIT_LIST_HEAD(&pcbFree_h);

    for (i = 0; i < MAXPROC; i++) {
        INIT_LIST_HEAD(&pcbFree_table[i].p_list);
        pcbFree_table[i].p_pid = 0;
        list_add_tail(&pcbFree_table[i].p_list, &pcbFree_h);
    }
}
\end{lstlisting}
Chiamata \emph{una sola volta} all'avvio del sistema (da \code{initial.c}).
Inizializza la sentinella \code{pcbFree\_h} come lista vuota, poi per
ciascuno dei 20 slot dell'array statico: inizializza il proprio
\code{p\_list} come nodo isolato e lo accoda (\code{list\_add\_tail}) alla
free list. Al termine, \code{pcbFree\_h} contiene tutti e 20 gli slot, nello
stesso ordine dell'array. Costo: $O(\text{MAXPROC})$, eseguito una volta
sola al boot.

\subsection{\func{allocPcb} --- allocazione}
\begin{lstlisting}
pcb_t* allocPcb() {
    if (list_empty(&pcbFree_h)) return NULL;

    struct list_head *pos = pcbFree_h.next;
    list_del(pos);

    pcb_t *new_pcb = container_of(pos, pcb_t, p_list);
    new_pcb->p_pid = next_pid++;
    INIT_LIST_HEAD(&new_pcb->p_list);
    INIT_LIST_HEAD(&new_pcb->p_child);
    INIT_LIST_HEAD(&new_pcb->p_sib);
    new_pcb->p_parent        = NULL;
    new_pcb->p_semAdd        = NULL;
    new_pcb->p_supportStruct = NULL;
    new_pcb->p_time          = 0;
    new_pcb->p_prio          = 0;

    return new_pcb;
}
\end{lstlisting}
Se la free list è vuota (tutti i 20 PCB sono in uso: limite massimo di
processi concorrenti raggiunto), ritorna \code{NULL} --- il chiamante (in
Fase 2, la syscall \code{CREATEPROCESS}) deve gestire questo fallimento
restituendo un errore al processo padre, \emph{senza mai} bloccarsi o andare
in crash. Altrimenti:
\begin{enumerate}
\item preleva il \emph{primo} nodo della free list (\code{pcbFree\_h.next})
  e lo stacca con \code{list\_del};
\item risale al \code{pcb\_t} contenitore con \code{container\_of};
\item gli assegna un nuovo PID univoco (\code{next\_pid++}, post-incremento:
  usa il valore corrente e \emph{poi} incrementa);
\item reinizializza da zero \emph{tutti} i campi rilevanti: le tre liste
  intrusive (\code{p\_list}, \code{p\_child}, \code{p\_sib}) come liste
  vuote, e azzera puntatore al padre, puntatore al semaforo, puntatore alla
  support struct, tempo CPU e priorità.
\end{enumerate}
Questo azzeramento sistematico è cruciale: un PCB può essere stato usato in
precedenza da un altro processo (poi terminato e liberato), quindi
conterrebbe ``spazzatura'' residua se non venisse ripulito esplicitamente ---
non essendoci \code{malloc}, non c'è nemmeno uno zeroing implicito come con
\code{calloc}.

\subsection{\func{freePcb} --- deallocazione}
\begin{lstlisting}
void freePcb(pcb_t* p) {
    if (p != NULL) list_add_tail(&p->p_list, &pcbFree_h);
}
\end{lstlisting}
Semplicemente riaggancia \code{p} in coda alla free list tramite il suo
campo \code{p\_list} (che a questo punto deve essere già stato staccato da
qualunque altra coda in cui si trovasse --- responsabilità del chiamante,
tipicamente \code{terminateProcess} in Fase 2). $O(1)$.

\subsection{Gestione delle code di processi (ready queue e code semaforo)}
\label{sub:code-processi-fase1}
\begin{lstlisting}
void mkEmptyProcQ(struct list_head* head) {
    INIT_LIST_HEAD(head);
}

int emptyProcQ(struct list_head* head) {
    return list_empty(head);
}

void insertProcQ(struct list_head* head, pcb_t* p) {
    struct list_head *pos;
    pcb_t *iter;
    list_for_each(pos, head) {
        iter = container_of(pos, pcb_t, p_list);
        if (p->p_prio > iter->p_prio) {
            __list_add(&p->p_list, pos->prev, pos);
            return;
        }
    }
    list_add_tail(&p->p_list, head);
}

pcb_t* headProcQ(struct list_head* head) {
    if (list_empty(head)) return NULL;
    return container_of(head->next, pcb_t, p_list);
}

pcb_t* removeProcQ(struct list_head* head) {
    if (list_empty(head)) return NULL;
    struct list_head *first = head->next;
    list_del(first);
    INIT_LIST_HEAD(first);
    return container_of(first, pcb_t, p_list);
}

pcb_t* outProcQ(struct list_head* head, pcb_t* p) {
    struct list_head *pos;
    pcb_t *iter;
    list_for_each(pos, head) {
        iter = container_of(pos, pcb_t, p_list);
        if (iter == p) {
            list_del(pos);
            return p;
        }
    }
    return NULL;
}
\end{lstlisting}

\begin{itemize}
\item \func{mkEmptyProcQ(head)}: banale wrapper su \code{INIT\_LIST\_HEAD};
  crea/azzera una coda di processi vuota. Usata da \code{initial.c} per la
  \code{readyQueue} globale.
\item \func{emptyProcQ(head)}: wrapper su \code{list\_empty}.
\item \func{insertProcQ(head, p)}: \textbf{questa è l'unica funzione della
  Fase 1 che introduce una politica di ordinamento}. Scorre la coda dal
  primo elemento e, appena trova un elemento \code{iter} con priorità
  \emph{strettamente minore} di \code{p}, inserisce \code{p} \emph{prima} di
  \code{iter} (con l'inserimento a basso livello \code{\_\_list\_add}, per
  poterlo fare in mezzo alla lista). Se non trova nessun elemento con
  priorità minore (cioè tutti hanno priorità $\ge$ quella di \code{p}, oppure
  la coda è vuota), lo accoda in fondo con \code{list\_add\_tail}.
  Risultato: la coda resta sempre ordinata per priorità
  \emph{decrescente} in testa, e a parità di priorità l'ordine di
  inserimento è preservato (FIFO) --- è esattamente ciò che serve per
  realizzare un \textbf{round-robin con priorità} nello scheduler di Fase 2:
  basta sempre estrarre la testa.
\item \func{headProcQ(head)}: restituisce (senza rimuoverlo) il PCB in
  testa, cioè il primo container\_of di \code{head->next}; \code{NULL} se
  la coda è vuota. Usata dallo scheduler per ispezionare il prossimo
  candidato senza consumarlo (es. nella gestione dello YIELD).
\item \func{removeProcQ(head)}: rimuove e restituisce la testa. Stacca il
  primo nodo (\code{list\_del}), lo reinizializza come lista a sé
  (\code{INIT\_LIST\_HEAD}, per igiene, così il PCB rimosso non conserva
  puntatori pendenti alla vecchia coda) e ne risale al \code{pcb\_t}.
  Questa è la funzione che lo scheduler chiama per estrarre il
  \emph{prossimo processo da eseguire}.
\item \func{outProcQ(head, p)}: rimozione \emph{mirata} di un PCB
  specifico da una coda arbitraria (non necessariamente in testa). Scorre
  linearmente la coda con \code{list\_for\_each} finché non trova il nodo il
  cui contenitore è \code{p}, poi lo stacca con \code{list\_del} e lo
  ritorna; se \code{p} non è nella coda, ritorna \code{NULL}. È $O(n)$ nel
  caso peggiore (a differenza delle altre operazioni, tutte $O(1)$): serve
  ad esempio quando un processo viene terminato mentre è ancora nella ready
  queue (in Fase 2, \code{terminateProcess} chiama sempre
  \code{outProcQ(\&readyQueue, p)} per sicurezza, anche se \code{p} potrebbe
  non essere in quella coda in quel momento --- la funzione in tal caso
  semplicemente non trova nulla e ritorna \code{NULL} senza effetti).
\end{itemize}

\subsection{Gestione dell'albero dei processi}
\begin{lstlisting}
int emptyChild(pcb_t* p) {
    return list_empty(&p->p_child);
}

void insertChild(pcb_t *prnt, pcb_t *p) {
    p->p_parent = prnt;
    list_add_tail(&p->p_sib, &prnt->p_child);
}

pcb_t* removeChild(pcb_t* p) {
    if (!p || list_empty(&p->p_child)) return NULL;
    struct list_head *first = p->p_child.next;
    pcb_t *child = container_of(first, pcb_t, p_sib);
    list_del(&child->p_sib);
    INIT_LIST_HEAD(&child->p_sib);
    child->p_parent = NULL;
    return child;
}

pcb_t* outChild(pcb_t* p) {
    if (!p || !p->p_parent) return NULL;
    list_del(&p->p_sib);
    INIT_LIST_HEAD(&p->p_sib);
    p->p_parent = NULL;
    return p;
}
\end{lstlisting}

La rappresentazione dell'albero dei processi è ``a puntatore al padre + lista
dei figli'', il classico schema \emph{first-child, next-sibling} adattato
alle liste intrusive:
\begin{itemize}
\item ogni PCB ha un \emph{unico} puntatore \code{p\_parent} al padre;
\item ogni PCB ha una \emph{lista} \code{p\_child}, che è la testa della
  lista dei propri figli;
\item ogni PCB si aggancia alla lista \code{p\_child} del padre tramite il
  proprio campo \code{p\_sib} (\emph{sibling}, fratello).
\end{itemize}

\begin{itemize}
\item \func{emptyChild(p)}: vero se \code{p} non ha figli.
\item \func{insertChild(prnt, p)}: rende \code{p} figlio di \code{prnt}:
  imposta \code{p->p\_parent = prnt} e accoda \code{p->p\_sib} in fondo a
  \code{prnt->p\_child}. Chiamata dalla syscall \code{CREATEPROCESS} appena
  dopo aver allocato il nuovo PCB.
\item \func{removeChild(p)}: rimuove e ritorna il \emph{primo} figlio di
  \code{p} (quello inserito per primo, in testa alla lista), staccandolo
  (\code{list\_del(\&child->p\_sib)}) e azzerandone il puntatore al padre.
  Se \code{p} non ha figli ritorna \code{NULL}. È usata da
  \code{terminateProcess} (Fase 2) in un ciclo
  \code{while ((child = removeChild(p)) != NULL) terminateProcess(child);}
  per terminare \emph{ricorsivamente e distruttivamente} tutto un
  sotto-albero: a ogni iterazione stacca un figlio (quindi la lista si
  accorcia) e lo termina ricorsivamente, il che a sua volta staccherà e
  terminerà i suoi discendenti, finché l'albero sotto \code{p} non è vuoto.
\item \func{outChild(p)}: operazione \emph{duale} di \code{removeChild}:
  data direttamente la vittima \code{p} (non il padre), la stacca dalla
  lista dei figli del \emph{proprio} padre e azzera \code{p\_parent}. Se
  \code{p} non ha un padre (\code{p\_parent == NULL}, es. è il processo
  radice, oppure è già stato staccato) ritorna \code{NULL} senza fare nulla.
  È usata in \code{terminateProcess} \emph{dopo} aver già staccato tutti i
  figli di \code{p} da \code{p} stesso, per staccare \code{p} dal proprio
  padre: se non lo si facesse, un processo che termina se stesso (o che
  viene terminato mentre non è più figlio ``attivo'') resterebbe agganciato
  come figlio ``fantasma'' nella lista del padre, e una successiva
  terminazione del padre lo ritroverebbe (già liberato!) tramite
  \code{removeChild}, causando un uso di memoria invalida.
\end{itemize}

% ============================================================
\section{\code{phase1/asl.c} --- Active Semaphore List}
% ============================================================

\subsection{Il file header \code{phase1/headers/asl.h}}
\begin{lstlisting}
// SPDX-FileCopyrightText: 2022 Luca Bassi, Gaia Clerici, Mirco Dondi, Fabio Gaiba
//
// SPDX-License-Identifier: GPL-3.0-or-later

#ifndef ASL_H_INCLUDED
#define ASL_H_INCLUDED

#include "../../headers/listx.h"
#include "../../headers/types.h"

extern semd_t semd_table[MAXPROC];
extern struct list_head semdFree_h;
extern struct list_head semd_h;

void initASL();
int insertBlocked(int* semAdd, pcb_t* p);
pcb_t* removeBlocked(int* semAdd);
pcb_t* outBlocked(pcb_t* p);
pcb_t* headBlocked(int* semAdd);

#endif
\end{lstlisting}
Stessa struttura di \code{pcb.h}, con una differenza degna di nota: qui
\code{semd\_table}, \code{semdFree\_h} e \code{semd\_h} sono dichiarate
\code{extern} (non \code{static} come i loro analoghi in \code{pcb.c}):
sono quindi \emph{tecnicamente} accessibili anche da fuori
\code{asl.c}, a differenza dei PCB. Nella pratica nessun altro modulo del
progetto le usa direttamente (Fase 2 e Fase 3 passano sempre attraverso le
5 funzioni pubbliche), ma la loro visibilità \code{extern} lascia aperta
la possibilità --- prevista dalla specifica originale di Pandos, da cui
questa interfaccia deriva --- di ispezionarle a fini di debug avanzato.
Anche questo header porta la stessa attribuzione SPDX del 2022 discussa
per \code{pcb.h}: fa parte della stessa interfaccia ``data'', che
\code{asl.c} implementa.

\subsection{Cos'è la ASL e perché serve}
Un semaforo, in questo progetto, è semplicemente un \code{int} (es. una cella
dell'array \code{devSems} in Fase 2, oppure una variabile globale come
\code{swapPoolSem} in Fase 3): non esiste una struttura dati ``semaforo''
associata a ciascuna variabile intera. Quando però un processo si blocca
\emph{su} un semaforo, serve un posto dove tenere la coda dei processi in
attesa \emph{per quel preciso indirizzo}. La \textbf{Active Semaphore List}
(ASL) è esattamente questo: una lista di descrittori \code{semd\_t}, uno per
ogni semaforo che ha \emph{almeno un} processo bloccato in questo momento
(da cui ``active''). Un semaforo che nessuno sta aspettando non ha nessun
descrittore nella ASL: se il suo valore diventa negativo per la prima volta,
un descrittore viene allocato al volo dalla free list; quando l'ultimo
processo in attesa su di esso viene rimosso, il descrittore torna nella free
list.

\subsection{Variabili globali del modulo}
\begin{lstlisting}
semd_t semd_table[MAXPROC];
struct list_head semdFree_h;
struct list_head semd_h;
\end{lstlisting}
\begin{itemize}
\item \code{semd\_table[MAXPROC]}: array statico di descrittori di
  semaforo. La dimensione è \code{MAXPROC} (20) e non un valore dedicato ai
  semafori perché, nella teoria classica di Pandos, nel caso peggiore
  \emph{ogni} processo può essere bloccato su un semaforo \emph{diverso}: non
  serve più di \code{MAXPROC} descrittori contemporaneamente attivi.
\item \code{semdFree\_h}: free list dei descrittori non in uso.
\item \code{semd\_h}: la ASL vera e propria, cioè la lista dei descrittori
  \emph{attualmente attivi} (con almeno un processo in coda).
\end{itemize}
A differenza di \code{pcb.c}, qui non c'è dichiarazione \code{static}: sono
dichiarate anche in \code{asl.h} con \code{extern}, quindi visibili da fuori
il modulo (anche se in pratica solo \code{asl.c} le manipola).

\subsection{\func{initASL}}
\begin{lstlisting}
void initASL() {
    INIT_LIST_HEAD(&semdFree_h);
    INIT_LIST_HEAD(&semd_h);
    for (int i = 0; i < MAXPROC; i++) {
        semd_table[i].s_key = NULL;
        INIT_LIST_HEAD(&semd_table[i].s_procq);
        list_add_tail(&semd_table[i].s_link, &semdFree_h);
    }
}
\end{lstlisting}
Analoga a \code{initPcbs}: inizializza le due liste sentinella come vuote,
poi per ciascuno dei 20 descrittori azzera la chiave, inizializza la coda dei
processi bloccati come vuota, e lo accoda alla free list tramite
\code{s\_link}. Chiamata una sola volta al boot, subito dopo
\code{initPcbs()}.

\subsection{\func{insertBlocked} --- blocco di un processo su un semaforo}
\begin{lstlisting}
int insertBlocked(int* semAdd, pcb_t* p) {
    if (!semAdd || !p) return 1; // errore
    semd_t* s;
    list_for_each_entry(s, &semd_h, s_link) {
        if (s->s_key == semAdd) {
            list_add_tail(&p->p_list, &s->s_procq);
            p->p_semAdd = semAdd;
            return 0;
        }
    }
    // semd non trovato, usa uno dalla free list
    if (list_empty(&semdFree_h)) return 1; // nessun semd disponibile
    s = container_of(semdFree_h.next, semd_t, s_link);
    list_del(&s->s_link);
    s->s_key = semAdd;
    INIT_LIST_HEAD(&s->s_procq);
    list_add_tail(&p->p_list, &s->s_procq);
    p->p_semAdd = semAdd;
    list_add_tail(&s->s_link, &semd_h);
    return 0;
}
\end{lstlisting}
Logica in due fasi:
\begin{enumerate}
\item \textbf{Ricerca}: scorre \code{semd\_h} con
  \code{list\_for\_each\_entry} confrontando \code{s->s\_key} con
  \code{semAdd}. Se trova un descrittore \emph{già attivo} per quel
  semaforo, ci accoda semplicemente \code{p} (tramite \code{p->p\_list}),
  imposta \code{p->p\_semAdd = semAdd} e ritorna \code{0} (successo).
\item \textbf{Allocazione}: se nessun descrittore esiste ancora per quel
  semaforo (è la \emph{prima} volta che qualcuno si blocca su di esso), ne
  preleva uno dalla free list (se disponibile, altrimenti errore \code{1}:
  situazione limite che segnalerebbe un bug, dato che \code{MAXPROC}
  descrittori bastano sempre a coprire \code{MAXPROC} processi), lo
  configura con la chiave giusta e una coda vuota, ci accoda \code{p}, e lo
  sposta dalla free list alla ASL attiva (\code{list\_add\_tail(\&s->s\_link,
  \&semd\_h)}).
\end{enumerate}
Nota: la ricerca è \emph{lineare}, $O(|\text{semd\_h}|)$: nel caso peggiore
$O(\text{MAXPROC})$. È accettabile perché il numero di semafori
\emph{attivi} contemporaneamente è comunque limitato dal numero di processi
vivi, e la specifica del progetto non richiede una struttura ordinata più
sofisticata (a differenza di alcune varianti storiche di Pandos che tengono
la ASL ordinata per indirizzo per poter interrompere la ricerca in anticipo).

\subsection{\func{removeBlocked} --- sblocco del primo processo in coda}
\begin{lstlisting}
pcb_t* removeBlocked(int* semAdd) {
    if (!semAdd) return NULL;
    semd_t* s;
    list_for_each_entry(s, &semd_h, s_link) {
        if (s->s_key == semAdd) {
            if (list_empty(&s->s_procq)) return NULL;
            pcb_t* p = container_of(s->s_procq.next, pcb_t, p_list);
            list_del(&p->p_list);
            p->p_semAdd = NULL;
            if (list_empty(&s->s_procq)) {
                list_del(&s->s_link);
                s->s_key = NULL;
                list_add_tail(&s->s_link, &semdFree_h);
            }
            return p;
        }
    }
    return NULL;
}
\end{lstlisting}
Trova il descrittore del semaforo \code{semAdd} (stessa ricerca lineare di
sopra). Se la sua coda è vuota (situazione anomala: un descrittore attivo
dovrebbe sempre avere almeno un processo, altrimenti sarebbe già tornato alla
free list) ritorna \code{NULL}. Altrimenti:
\begin{enumerate}
\item preleva il \emph{primo} PCB della coda (FIFO: chi si è bloccato per
  primo viene sbloccato per primo), lo stacca e azzera il suo
  \code{p\_semAdd} (segnalando che ora non è più bloccato su nulla);
\item \textbf{se dopo la rimozione la coda del semaforo è diventata vuota},
  il descrittore non serve più: viene staccato dalla ASL e restituito alla
  free list. Questo è il meccanismo che mantiene la ASL ``pulita'',
  contenente solo descrittori realmente attivi.
\end{enumerate}
Chiamata dalla syscall \code{VERHOGEN} (V) in Fase 2, quando il valore del
semaforo, dopo l'incremento, resta $\le 0$ (cioè c'era almeno un
processo in attesa).

\subsection{\func{outBlocked} --- rimozione mirata (non necessariamente in testa)}
\begin{lstlisting}
pcb_t* outBlocked(pcb_t* p) {
    if (!p || !p->p_semAdd) return NULL;
    semd_t* s;
    list_for_each_entry(s, &semd_h, s_link) {
        if (s->s_key == p->p_semAdd) {
            struct list_head* pos;
            list_for_each(pos, &s->s_procq) {
                if (pos == &p->p_list) {
                    list_del(&p->p_list);
                    p->p_semAdd = NULL;
                    if (list_empty(&s->s_procq)) {
                        list_del(&s->s_link);
                        s->s_key = NULL;
                        list_add_tail(&s->s_link, &semdFree_h);
                    }
                    return p;
                }
            }
            return NULL;
        }
    }
    return NULL;
}
\end{lstlisting}
Rimuove \code{p} dalla coda del semaforo su cui è bloccato (letto da
\code{p->p\_semAdd}, non passato esplicitamente: è \code{p} stesso a
``sapere'' su cosa è bloccato), \emph{indipendentemente dalla sua posizione}
nella coda (non necessariamente il primo). Doppio ciclo annidato: prima trova
il descrittore giusto (uguale a \code{insertBlocked}/\code{removeBlocked}),
poi scorre la sua coda di processi confrontando i nodi \code{p\_list} finché
non trova esattamente \code{\&p->p\_list}. Anche qui, se la coda si svuota
dopo la rimozione, il descrittore torna alla free list.

È l'unica delle tre operazioni di rimozione a costo $O(n \cdot m)$ nel caso
peggiore (dove $n$ = semafori attivi, $m$ = lunghezza della coda), ma serve
per un caso che le altre due non coprono: quando un processo bloccato
\emph{da qualche parte} deve essere estratto senza sapere a priori la sua
posizione. È usata in Fase 2 da \code{terminateProcess}: se il processo che
si sta terminando risulta bloccato su un semaforo (\code{p->p\_semAdd !=
NULL}), bisogna toglierlo dalla coda di quel semaforo prima di liberarne il
PCB, altrimenti resterebbe un puntatore pendente nella ASL verso un PCB ormai
restituito alla free list.

\subsection{\func{headBlocked} --- ispezione senza rimozione}
\begin{lstlisting}
pcb_t* headBlocked(int* semAdd) {
    if (!semAdd) return NULL;
    semd_t* s;
    list_for_each_entry(s, &semd_h, s_link) {
        if (s->s_key == semAdd) {
            if (list_empty(&s->s_procq)) return NULL;
            return container_of(s->s_procq.next, pcb_t, p_list);
        }
    }
    return NULL;
}
\end{lstlisting}
Come \code{removeBlocked} ma \emph{senza} modificare nulla: trova il
descrittore del semaforo e restituisce (senza toglierlo) il primo PCB della
sua coda, oppure \code{NULL} se il semaforo non è attivo o la sua coda è
vuota. Nel codice attuale del progetto non è invocata direttamente dai
moduli di Fase 2/3 forniti (serve come primitiva di ispezione prevista dalla
specifica standard di Pandos), ma fa parte dell'interfaccia pubblica di
\code{asl.h} richiesta.

% ============================================================
\section{\code{klog.c} --- buffer di log circolare per il debug}
% ============================================================

\code{klog.c} è compilato in \textbf{tutti e tre} i target del progetto
(\code{CMakeLists.txt}: compare come primo file sorgente di
\code{MultiPandOS\_phase1}, \code{MultiPandOS\_phase2} e
\code{MultiPandOS\_phase3}), quindi va conosciuto anche se non appartiene
concettualmente a una fase specifica.

\begin{lstlisting}
#define KLOG_LINES     64
#define KLOG_LINE_SIZE 42

unsigned int klog_line_index = 0;
unsigned int klog_char_index = 0;
char         klog_buffer[KLOG_LINES][KLOG_LINE_SIZE] = {0};

void klog_print(char *str) {
    while (*str != '\0') {
        if (*str == '\n') { next_line(); str++; }
        else { klog_buffer[klog_line_index][klog_char_index] = *str++; next_char(); }
    }
}

void klog_print_dec(unsigned int num) { /* stampa un numero decimale 0..99 */ }
void klog_print_hex(unsigned int num) { /* stampa un numero in esadecimale */ }

static void next_char(void) {
    if (++klog_char_index >= KLOG_LINE_SIZE) { klog_char_index = 0; next_line(); }
}
static void next_line(void) {
    klog_line_index = (klog_line_index + 1) % KLOG_LINES;
    klog_char_index = 0;
    for (unsigned int i = 0; i < KLOG_LINE_SIZE; i++)
        klog_buffer[klog_line_index][i] = ' ';
}
\end{lstlisting}

\`E una piccola libreria di logging \textbf{indipendente dal terminale
virtuale}: invece di scrivere caratteri su un device (che richiede
\code{DOIO}, quindi un Nucleus già funzionante e un processo che può
bloccarsi), \code{klog\_print} scrive direttamente in un array globale
circolare \code{klog\_buffer} (64 righe da 42 caratteri). Il buffer è
pensato per essere \textbf{ispezionato dal debugger dell'emulatore}
(tracciando l'indirizzo del simbolo \code{klog\_buffer} nella finestra di
memoria di uMPS3/uRISCV): permette di lasciare ``tracce'' di cosa sta
succedendo nel kernel anche in situazioni in cui il terminale non è ancora
disponibile o si sospetta che il problema sia proprio nel sottosistema di
I/O che si starebbe usando per stampare a video.
\begin{itemize}
\item \code{klog\_print(str)}: copia \code{str} carattere per carattere nella
  riga corrente, andando a capo (\code{next\_line}) su ogni
  \code{'\textbackslash n'} incontrato nella stringa sorgente.
\item \code{klog\_print\_dec}/\code{klog\_print\_hex}: varianti per stampare
  un numero, rispettivamente in decimale (limitato per costruzione
  all'intervallo 0--99: il ramo \code{if(num >= 10)} gestisce due o più
  cifre in un ciclo, il ramo \code{else} gestisce esplicitamente il caso a
  una sola cifra) o in esadecimale (nessun limite, cifre \code{0-9A-F}).
\item \code{next\_char}/\code{next\_line} (\code{static}, private del
  modulo): gestiscono l'avanzamento circolare dentro la riga corrente e tra
  le righe (\code{\% KLOG\_LINES}), ripulendo con spazi la riga nuova per
  leggibilità quando si va a capo.
\end{itemize}

\paragraph{Attenzione: nel codice attuale non è mai invocato.} Una ricerca
di \code{klog\_print} nell'intero albero dei sorgenti (\code{pcb.c},
\code{asl.c}, \code{initial.c}, \code{scheduler.c}, \code{exceptions.c},
\code{interrupts.c}, \code{initProc.c}, \code{vmSupport.c},
\code{sysSupport.c}) non produce alcun risultato: è compilato e linkato in
ogni build, ma \textbf{nessuna chiamata a queste funzioni è presente} nel
resto del kernel fornito. Il debug ``a schermo'' effettivamente utilizzato
nel progetto passa invece dalle macro condizionali di
\code{phase2/debug.h} (si veda il documento di Fase 2, che scrivono davvero
sul terminale via i registri hardware). \`E
comunque utile sapere spiegare \emph{cosa fa e a cosa servirebbe}
\code{klog.c} se l'esaminatore lo notasse tra i sorgenti compilati.

% ============================================================
\section{Relazioni tra le strutture: la ``mappa'' d'insieme}
% ============================================================

\`E utile fissare mentalmente come tutte queste liste si intreccino durante
la vita di un processo:

\begin{enumerate}
\item Al boot, tutti i 20 PCB sono nella \textbf{free list} dei PCB
  (\code{pcbFree\_h}) e tutti i 20 \code{semd\_t} sono nella \textbf{free
  list} dei semafori (\code{semdFree\_h}).
\item \code{allocPcb()} sposta un PCB dalla free list nel ``mondo'': da quel
  momento il suo campo \code{p\_list} è libero per essere agganciato altrove
  (Fase 2 lo inserisce subito nella ready queue).
\item Mentre un processo è \textbf{pronto}, il suo \code{p\_list} è nella
  \code{readyQueue} (ordinata per priorità da \code{insertProcQ}).
\item Quando un processo esegue una \code{P} che lo blocca, \code{p\_list}
  viene spostato (da \code{insertBlocked}) nella coda \code{s\_procq} del
  descrittore \code{semd\_t} giusto, allocandolo dalla free list dei
  semafori se necessario.
\item Una \code{V} (\code{removeBlocked}) o una terminazione forzata
  (\code{outBlocked}) rimuove il processo da quella coda e, se la coda
  diventa vuota, il \code{semd\_t} torna alla free list dei semafori.
\item Parallelamente, ogni processo (tranne il primo) è agganciato tramite
  \code{p\_sib} alla lista \code{p\_child} del proprio padre: questo albero
  \emph{coesiste} con le code sopra, è del tutto indipendente da esse (usa
  campi diversi) e serve solo a sapere ``chi è figlio di chi'', per poter
  terminare interi sotto-alberi.
\item Alla terminazione, il PCB viene staccato da tutte le liste in cui si
  trova (\code{outChild} dal padre, \code{outBlocked}/\code{outProcQ} dalla
  coda in cui si trovava) e restituito, tramite \code{freePcb}, alla free
  list dei PCB, pronto per essere riassegnato a un futuro processo.
\end{enumerate}

% ============================================================
\section{Il test driver ufficiale: \code{phase1/p1test.c}}
% ============================================================

\code{p1test.c} \textbf{non è codice originale del progetto}: è il test
driver standard fornito dal corso per verificare \code{pcb.c} e
\code{asl.c} (compilato solo nel target \code{phase1} di
\code{CMakeLists.txt}, insieme a \code{klog.c}, \code{pcb.c},
\code{asl.c}). Vale comunque la pena sapere \emph{come è strutturato} e
\emph{cosa dimostra}, perché è il modo in cui si osserva la Fase 1
``funzionare'' a schermo (gira da sola, senza Nucleus: non ci sono
eccezioni né scheduler, è pura esecuzione sequenziale in \code{main}).

\subsection{Meccanismo di segnalazione}
Due funzioni di supporto, \code{addokbuf}/\code{adderrbuf}, accumulano i
messaggi di progresso (rispettivamente di errore) in due buffer statici
(\code{okbuf}, \code{errbuf}) e li stampano anche a video sul terminale 0
tramite \code{termprint} (una routine che scrive un carattere alla volta
sul registro di comando del device, in polling attivo --- non usa
\code{DOIO}, dato che il Nucleus non esiste ancora in questa fase).
\code{adderrbuf} segnala l'errore e poi entra in un \textbf{loop infinito}
(\code{while(1);}): non essendoci un vero kernel capace di gestire un
\code{PANIC}, è il modo più semplice per ``fermare tutto'' non appena
qualcosa non torna, così che l'errore resti visibile a video.

\subsection{Cosa verifica, in ordine}
\begin{enumerate}
\item \textbf{\code{allocPcb}/\code{freePcb}}: alloca tutti i
  \code{MAXPROC} PCB disponibili verificando che nessuna allocazione
  fallisca, poi verifica che un'allocazione \emph{oltre} il limite
  restituisca \code{NULL}; libera 10 PCB per riportarli in free list.
\item \textbf{Code di processi} (\code{insertProcQ}/\code{removeProcQ}/\code{
  headProcQ}/\code{outProcQ}/\code{emptyProcQ}): costruisce una coda di 10
  elementi, verifica che \code{headProcQ} indichi il primo inserito, toglie
  un elemento dalla testa e uno dal centro con \code{outProcQ} (e verifica
  che togliere un elemento \emph{non presente} restituisca \code{NULL}),
  poi svuota la coda con \code{removeProcQ} controllando che l'ordine di
  uscita e l'ultimo elemento estratto siano quelli attesi.
\item \textbf{Ordinamento per priorità}: inserisce quattro PCB con priorità
  \code{20, 10, 5, 10} (l'ultimo, \code{p2bis}, uguale al secondo) e
  verifica che \code{removeProcQ} li restituisca nell'ordine
  \code{p1(20)}, \code{p2(10)}, \code{p2bis(10)}, \code{p3(5)}: dimostra sia
  l'ordinamento per priorità decrescente sia il rispetto del FIFO a parità
  di priorità (esattamente il comportamento di \code{insertProcQ}
  analizzato nella Sezione~\ref{sub:code-processi-fase1}).
\item \textbf{Albero dei processi} (\code{insertChild}/\code{removeChild}/\code{
  outChild}/\code{emptyChild}): rende 9 PCB figli di uno stesso padre,
  verifica \code{outChild} su un figlio in mezzo alla lista e su uno non
  presente, poi svuota l'albero con \code{removeChild} controllando che,
  esaurita la lista, un'ulteriore chiamata restituisca \code{NULL}.
\item \textbf{ASL} (\code{initASL}, \code{insertBlocked},
  \code{removeBlocked}, \code{headBlocked}, \code{outBlocked}): blocca
  processi su diversi semafori, verifica che un descrittore
  \code{semd\_t} torni davvero alla free list quando la sua coda si
  svuota (rispolverandolo subito dopo con un nuovo \code{insertBlocked}
  sullo stesso indirizzo), e che \code{outBlocked} rimuova correttamente un
  processo specifico (anche non in testa) senza poterlo rimuovere una
  seconda volta.
\end{enumerate}
Se tutti i controlli passano, l'ultimo messaggio stampato è la celebre
citazione \emph{``So Long and Thanks for All the Fish''} (Guida galattica
per gli autostoppisti): è il segnale, familiare a chi ha seguito corsi
analoghi altrove, che l'intera Fase 1 supera i test.

\paragraph{Nota su \code{test\_asl\_manuale.c}.} Nella cartella
\code{phase1/} è presente anche un piccolo programma
\code{test\_asl\_manuale.c}, scritto per un test manuale rapido della sola
ASL (usa \code{printf} e compila per l'host, \emph{non} fa parte di nessun
target CMake): non è codice di produzione né viene eseguito
sull'emulatore, ma uno strumento di debug estemporaneo usato durante lo
sviluppo per ispezionare a occhio lo stato delle code dei semafori dopo
\code{insertBlocked}/\code{removeBlocked}/\code{outBlocked}.

% ============================================================
\section{Sistema di build: \code{CMakeLists.txt} e \code{config\_machine.json}}
% ============================================================

Anche se non è ``codice del kernel'', il sistema di build va conosciuto:
determina esattamente quali file finiscono in quale eseguibile e quale
configurazione dell'emulatore lo carica --- è tipico che un esaminatore
chieda ``come fai a passare dal sorgente all'esecuzione in Fase 1?''.

\subsection{\code{CMakeLists.txt} --- il target \code{phase1}}
\begin{lstlisting}
add_link_options(-nostdlib -T ${URISCV_SRC}/uriscvcore.ldscript -march=rv32imfd)
add_compile_options(
    -ffreestanding -static -nostartfiles -nostdlib
    -I${URISCV_INC}
    -ggdb -Wall -O0 -std=gnu99
    -march=rv32imafd -mabi=ilp32d
)

add_executable(MultiPandOS_phase1 EXCLUDE_FROM_ALL
    ./klog.c
    phase1/pcb.c
    phase1/asl.c
    phase1/p1test.c
    ${URISCV_SRC}/crtso.S
    ${URISCV_SRC}/liburiscv.S
)

add_custom_target(phase1
    COMMAND ${CMAKE_COMMAND} -E copy
        ${PROJECT_BINARY_DIR}/MultiPandOS_phase1
        ${PROJECT_BINARY_DIR}/MultiPandOS
    COMMAND uriscv-elf2uriscv -k ${PROJECT_BINARY_DIR}/MultiPandOS
    BYPRODUCTS MultiPandOS.core.uriscv MultiPandOS.stab.uriscv
    DEPENDS MultiPandOS_phase1
)
\end{lstlisting}
Le opzioni di compilazione globali (\code{add\_compile\_options},
\code{add\_link\_options}, applicate a tutti e tre i target) riflettono la
natura ``freestanding'' del kernel: \code{-ffreestanding -nostdlib
-nostartfiles} escludono la libc standard e il normale runtime C (non
esiste un sistema operativo sottostante che fornisca queste cose: \emph{è}
il progetto stesso a doverle sostituire, es. con \code{crtso.S} al posto
del solito \code{crt0}), \code{-march=rv32imafd -mabi=ilp32d} fissano
l'architettura RISC-V a 32 bit con estensioni intero/moltiplicazione/
float, \code{-O0} disabilita le ottimizzazioni (più facile debuggare passo
passo nell'emulatore) e il linker script \code{uriscvcore.ldscript}
posiziona sezioni e simboli secondo il formato \code{.core} atteso dalla
macchina virtuale.

Il target \code{phase1} compila insieme \code{klog.c}, \code{pcb.c},
\code{asl.c} e \code{p1test.c} (più il runtime di sistema
\code{crtso.S}/\code{liburiscv.S}, forniti dall'installazione di uRISCV,
non dal progetto) in un unico eseguibile \code{MultiPandOS\_phase1}; il
target \emph{custom} \code{phase1} lo copia come \code{MultiPandOS}, lo
converte con \code{uriscv-elf2uriscv -k} nel formato \code{.core} atteso
dall'emulatore (producendo \code{build/MultiPandOS.core.uriscv} e la
relativa tabella dei simboli \code{.stab.uriscv}, usata dal debugger per
mostrare nomi di funzione/variabile invece dei soli indirizzi). I target
\code{phase2} e \code{phase3} (analizzati nei rispettivi documenti) seguono
esattamente la stessa struttura, cambiando solo l'elenco dei sorgenti.
\code{EXCLUDE\_FROM\_ALL} evita che un semplice \code{make} (senza
specificare il target) provi a compilare tutti e tre insieme, dato che
producono lo stesso file di output e si sovrascriverebbero a vicenda.

\subsection{\code{config\_machine.json} --- la configurazione dell'emulatore}
\begin{lstlisting}
{
    "boot": {
        "core-file": "build/MultiPandOS.core.uriscv",
        "load-core-file": true
    },
    "bootstrap-rom": "/usr/local/share/uriscv/coreboot.rom.uriscv",
    "clock-rate": 1,
    "devices": {
        "terminal0": {
            "enabled": true,
            "file": "term0.uriscv"
        }
    },
    "execution-rom": "/usr/local/share/uriscv/exec.rom.uriscv",
    "num-processors": 1,
    "num-ram-frames": 64,
    "symbol-table": {
        "asid": 64,
        "file": "build/MultiPandOS.stab.uriscv"
    },
    "tlb-floor-address": "0x0",
    "tlb-size": 16
}
\end{lstlisting}
\`E il file che l'interfaccia grafica di uRISCV/uMPS3 carica per sapere
\emph{come} istanziare la macchina virtuale, usato sia per la Fase 1 sia
per la Fase 2 (la Fase 3 usa invece \code{phase3\_config\_machine.json},
descritto nel documento di Fase 3). Campi rilevanti:
\begin{itemize}
\item \code{boot.core-file}: il file \code{.core} da caricare all'avvio
  --- esattamente il file prodotto dal target CMake attivo per ultimo
  (da qui l'avvertenza pratica: bisogna \emph{sempre} ricompilare il target
  giusto prima di avviare l'emulatore, altrimenti gira il kernel di
  un'altra fase).
\item \code{bootstrap-rom}/\code{execution-rom}: le ROM del firmware
  dell'emulatore stesso (BIOS), fornite dall'installazione di uRISCV, non
  dal progetto: sono il ``livello 1'' (firmware) della gerarchia descritta
  nell'introduzione.
\item \code{devices.terminal0}: l'unico device configurato in questa fase
  è un terminale (\code{term0.uriscv}, un file che l'emulatore usa per
  registrare/riprodurre l'input-output del terminale). Niente device flash:
  non serve memoria virtuale né backing store finché non si arriva alla
  Fase 3.
\item \code{num-ram-frames}: 64 frame di RAM fisica disponibili (sufficienti
  per il piccolo kernel di Fase 1/2, che non ha bisogno dei 32 frame
  riservati al Nucleus più i 16 dello Swap Pool della Fase 3 --- si
  confronti con i 256 frame di \code{phase3\_config\_machine.json}).
\item \code{symbol-table}: la tabella dei simboli generata insieme al file
  \code{.core} (\code{MultiPandOS.stab.uriscv}), usata dal debugger
  integrato per mostrare nomi simbolici.
\item \code{tlb-floor-address: "0x0"}: in Fase 1/2 non esiste alcuno spazio
  utente separato da proteggere con la paginazione, quindi il TLB non ha
  bisogno di un indirizzo di soglia diverso da zero (si confronti con
  \code{"0x80000000"} in Fase 3, che corrisponde proprio all'inizio di
  \code{kuseg}).
\end{itemize}

\end{document}
